\documentclass[12pt,a4paper]{article}
\usepackage[utf8]{inputenc}
\usepackage[T1]{fontenc}
\usepackage{amsmath,amssymb}
\usepackage[margin=1in]{geometry}
\begin{document}
\title{Kuliah 10 - Supervised dan Unsupervised Learning}
\author{Yeni Herdiyeni}
\date{}
\maketitle

\maketitle
\tableofcontents
\newpage

\section{Kuliah 10 - Supervised dan Unsupervised Learning}

\textbf{Topik:} Studi Kasus Analisis dan Prediksi Produktivitas Padi  

\textbf{Pengampu:} Yeni Herdiyeni, Program Studi Kecerdasan Buatan, Sekolah Sains Data, Matematika dan Informatika, IPB University  

\textbf{Tanggal:} 29 April 2026


\hrule


\subsection{1. Tujuan Pembelajaran}

Setelah mengikuti materi ini, mahasiswa diharapkan mampu:


\begin{enumerate}
\item Membedakan \textbf{supervised}, \textbf{unsupervised}, dan \textbf{reinforcement learning}.
\end{enumerate}

\begin{enumerate}
\item Menjelaskan peran \textbf{EDA} dalam pipeline machine learning.
\end{enumerate}

\begin{enumerate}
\item Melakukan \textbf{feature engineering} berdasarkan temuan EDA.
\end{enumerate}

\begin{enumerate}
\item Memahami konsep evaluasi model supervised dan unsupervised.
\end{enumerate}

\begin{enumerate}
\item Menginterpretasikan hasil model dalam konteks kasus pertanian.
\end{enumerate}


\hrule


\subsection{2. Use Case: Produktivitas Padi}

\subsubsection{2.1 Masalah Utama}

Petani dan pengambil kebijakan ingin memahami faktor-faktor yang memengaruhi hasil panen padi. Dengan memahami faktor-faktor tersebut, mereka dapat:



\subsubsection{2.2 Data yang Tersedia}

\begin{tabular}{|l|l|}
\hline
Variabel & Keterangan \\
\hline
Curah hujan & Jumlah hujan dalam periode tertentu (mm) \\
\hline
Suhu & Suhu rata-rata lingkungan ( derajat C) \\
\hline
Kelembaban & Kelembaban udara/tanah (%) \\
\hline
Jumlah pupuk & Dosis pupuk yang diberikan (kg/ha) \\
\hline
Luas lahan & Luas area tanam (ha) \\
\hline
Hasil panen & Target: produksi padi (ton/ha) \\
\hline
\end{tabular}


\subsubsection{2.3 Pertanyaan Data Science}

\begin{quote}
Bagaimana data ini dapat digunakan untuk memprediksi hasil panen dan memahami pola lahan?\end{quote}


\hrule


\subsection{3. Motivasi Machine Learning}

\subsubsection{3.1 Pendekatan Tradisional}

Dalam pendekatan tradisional, kita membuat aturan secara manual:


\begin{quote}
"Jika curah hujan tinggi dan pupuk cukup, maka hasil panen tinggi."\end{quote}


Namun, dunia nyata jauh lebih kompleks:


\subsubsection{3.2 Motivasi ML}

\begin{quote}
\textbf{Machine learning digunakan ketika pola dalam data sulit dirumuskan secara eksplisit dengan aturan manual.}\end{quote}


ML mampu:


\hrule


\subsection{4. Tiga Paradigma Machine Learning}

\begin{tabular}{|l|l|l|}
\hline
Paradigma & Data & Contoh dalam Use Case \\
\hline
\textbf{Supervised Learning} & Ada input dan target & Prediksi hasil panen berdasarkan cuaca dan pupuk \\
\hline
\textbf{Unsupervised Learning} & Hanya input, tanpa target & Segmentasi lahan berdasarkan karakteristiknya \\
\hline
\textbf{Reinforcement Learning} & Belajar dari aksi dan reward & Sistem irigasi otomatis belajar kapan memberi air \\
\hline
\end{tabular}


\hrule


\subsection{5. Mengapa Tidak Langsung Modeling?}

Kesalahan umum dalam machine learning adalah langsung membuat model tanpa memahami data. Risikonya:



\begin{quote}
\textbf{Pertanyaan kritis:} Sebelum membuat model, apakah kita sudah memahami karakteristik datanya?\end{quote}


\hrule


\subsection{6. Peran EDA dalam Machine Learning}

\subsubsection{6.1 Apa itu EDA?}

\textbf{Exploratory Data Analysis (EDA)} adalah tahap untuk memahami data sebelum modeling. EDA membantu menjawab pertanyaan seperti:



\begin{quote}
\textbf{Ide utama:} EDA bukan sekadar membuat grafik, tetapi proses berpikir untuk menentukan strategi modeling.\end{quote}


\subsubsection{6.2 EDA 1: Distribusi Curah Hujan}

\textbf{Pertanyaan:}


\textbf{Insight:}


\textbf{Keputusan:}


\subsubsection{6.3 EDA 2: Relasi Curah Hujan dan Hasil Panen}

\textbf{Pertanyaan:}


\textbf{Insight:}


\textbf{Keputusan:}


\subsubsection{6.4 EDA 3: Outlier pada Hasil Panen}

\textbf{Pertanyaan:}


\textbf{Insight:}


\textbf{Keputusan:}


\subsubsection{6.5 EDA 4: Korelasi Antar Variabel}

\textbf{Pertanyaan:}


\textbf{Insight:}


\textbf{Keputusan:}


\subsubsection{6.6 Dari EDA ke Keputusan Modeling}

\begin{tabular}{|l|l|l|}
\hline
Temuan EDA & Keputusan & Dampak ke Model \\
\hline
Data skewed & Scaling / transformasi & Model lebih stabil \\
\hline
Relasi non-linear & Tambah fitur kuadrat & Pola lebih tertangkap \\
\hline
Outlier bermakna & Tidak langsung dihapus & Model tetap merekam kasus ekstrem \\
\hline
Fitur lemah & Feature selection & Model lebih sederhana \\
\hline
\end{tabular}


\hrule


\subsection{7. Feature Engineering}

\subsubsection{7.1 Definisi}

\textbf{Feature engineering} adalah proses membentuk fitur agar lebih representatif bagi model.


\begin{quote}
\textbf{Intuisi:} Feature engineering adalah cara memasukkan pemahaman domain ke dalam data.\end{quote}


\subsubsection{7.2 Contoh pada Use Case Produktivitas Padi}

\textbf{Fitur kuadrat untuk menangkap relasi non-linear:}


\\[
\text{rainfall}^2
\\]


\textbf{Fitur rasio untuk menggambarkan efisiensi:}


\\[
\text{fertilizer per area} = \frac{\text{fertilizer}}{\text{land area}}
\\]


\textbf{Fitur kategori curah hujan:}


\begin{lstlisting}[language=Python]
df["rainfall_category"] = pd.cut(
    df["rainfall"],
    bins=[-np.inf, 140, 220, np.inf],
    labels=["rendah", "sedang", "tinggi"]
)
\end{lstlisting}

\hrule


\subsection{8. Supervised Learning: Prediksi Hasil Panen}

\subsubsection{8.1 Konsep}

Dalam supervised learning, data memiliki pasangan input dan target:


\\[
(X, y)
\\]


Pada use case ini:


\\[
X = \{\text{curah hujan}, \text{suhu}, \text{kelembaban}, \text{pupuk}, \text{luas lahan}\}
\\]


\\[
y = \text{hasil panen}
\\]


\textbf{Tujuan:} Model belajar fungsi $f(X) \rightarrow y$ untuk memprediksi hasil panen dari data baru.


\subsubsection{8.2 Contoh Model: Linear Regression}

Model sederhana:


\\[
\hat{y} = \beta_0 + \beta_1 x_1 + \beta_2 x_2 + \dots + \beta_p x_p
\\]


Karena EDA menunjukkan relasi non-linear, kita tambahkan fitur kuadrat:


\\[
\hat{y} = \beta_0 + \beta_1 \cdot \text{rainfall} + \beta_2 \cdot \text{rainfall}^2 + \dots
\\]


\begin{quote}
\textbf{Makna penting:} Model tetap linear terhadap parameter, tetapi mampu menangkap pola non-linear melalui transformasi fitur.\end{quote}


\subsubsection{8.3 Evaluasi Supervised Learning}

\textbf{Mean Absolute Error (MAE):}


\\[
\begin{tabular}{|l|l|}
\hline
\text{MAE} = \frac{1}{n} \sum_{i=1}^{n} & y_i - \hat{y}_i \\
\hline
\end{tabular}
\\]


\textbf{Root Mean Squared Error (RMSE):}


\\[
\text{RMSE} = \sqrt{\frac{1}{n} \sum_{i=1}^{n} (y_i - \hat{y}_i)^2}
\\]


\textbf{Koefisien Determinasi ($R^2$):}


\\[
R^2 = 1 - \frac{SS_{\text{res}}}{SS_{\text{tot}}}
\\]


\textbf{Interpretasi:} Semakin kecil MAE/RMSE dan semakin besar $R^2$, semakin baik model menjelaskan data.


\subsubsection{8.4 Grafik Actual vs Predicted}

Cara membaca grafik:


\hrule


\subsection{9. Unsupervised Learning: Segmentasi Lahan}

\subsubsection{9.1 Konsep}

Dalam unsupervised learning, data tidak memiliki target:


\\[
X = \{\text{curah hujan}, \text{pupuk}, \text{suhu}, \text{hasil panen}\}
\\]


\textbf{Tujuan:}


\textbf{Contoh pertanyaan:} Apakah terdapat kelompok lahan dengan kondisi input dan hasil panen yang berbeda?


\subsubsection{9.2 Contoh Model: K-Means Clustering}

K-Means membagi data ke dalam $K$ cluster dengan meminimalkan jarak intra-cluster:


\\[
\begin{tabular}{|l|l|l|}
\hline
\min \sum_{k=1}^{K} \sum_{x_i \in C_k} \ & x_i - \mu_k\ & ^2 \\
\hline
\end{tabular}
\\]


Di mana:


\textbf{Intuisi:} Data yang mirip ditempatkan dalam cluster yang sama. Setiap cluster memiliki pusat yang disebut centroid.


\begin{quote}
\textbf{Catatan:} Hasil clustering harus diinterpretasikan dengan \emph{domain knowledge}.\end{quote}


\subsubsection{9.3 Evaluasi Unsupervised: Silhouette Score}

Karena tidak ada label benar, evaluasi clustering menggunakan metrik seperti \textbf{Silhouette Score}:


\\[
s(i) = \frac{b(i) - a(i)}{\max(a(i), b(i))}
\\]


Di mana:


\textbf{Interpretasi:} Nilai mendekati 1 menunjukkan cluster lebih terpisah dengan baik.


\subsubsection{9.4 Elbow Method untuk Memilih Jumlah Cluster}

\textbf{Ide dasar:} Elbow method digunakan untuk menentukan jumlah cluster optimal $K$.


\textbf{Langkah:}


\\[
\begin{tabular}{|l|l|l|}
\hline
\text{WCSS} = \sum_{k=1}^{K} \sum_{x_i \in C_k} \ & x_i - \mu_k\ & ^2 \\
\hline
\end{tabular}
\\]


\begin{enumerate}
\item Ulangi untuk berbagai nilai $K$.
\end{enumerate}

\begin{enumerate}
\item Plot $K$ vs WCSS.
\end{enumerate}

\begin{enumerate}
\item Pilih titik "siku" (\emph{elbow}) di mana penurunan WCSS mulai melambat.
\end{enumerate}


\begin{quote}
\textbf{Pesan:} Pemilihan cluster perlu menggabungkan metrik dan interpretasi domain.\end{quote}


\subsubsection{9.5 Interpretasi Cluster}

Pertanyaan interpretatif:


\begin{quote}
\textbf{Insight:} Clustering bukan akhir analisis, tetapi awal untuk memahami segmentasi data.\end{quote}


\hrule


\subsection{10. Reinforcement Learning dalam Use Case}

Walaupun fokus praktikum adalah supervised dan unsupervised, reinforcement learning juga dapat muncul dalam pertanian.


\textbf{Contoh: Sistem Irigasi Otomatis}


\begin{tabular}{|l|l|}
\hline
Komponen & Peran \\
\hline
\textbf{Agent} & Sistem irigasi \\
\hline
\textbf{Environment} & Lahan pertanian \\
\hline
\textbf{Action} & Memberi air atau tidak \\
\hline
\textbf{Reward} & Peningkatan hasil panen atau efisiensi air \\
\hline
\end{tabular}


\begin{quote}
\textbf{Intuisi:} Reinforcement learning cocok ketika sistem harus belajar mengambil keputusan secara bertahap melalui feedback.\end{quote}


\hrule


\subsection{11. Pipeline Machine Learning End-to-End}


\begin{enumerate}
\item \textbf{Data Collection}
\end{enumerate}


\begin{enumerate}
\item \textbf{EDA}
\end{enumerate}


\begin{enumerate}
\item \textbf{Feature Engineering}
\end{enumerate}


\begin{enumerate}
\item \textbf{Modeling}
\end{enumerate}


\begin{enumerate}
\item \textbf{Evaluation}
\end{enumerate}


\begin{enumerate}
\item \textbf{Deployment}
\end{enumerate}


\hrule


\subsection{12. Deployment: Dari Model ke Keputusan}

Model machine learning tidak berhenti pada akurasi. Contoh deployment:



\begin{quote}
\textbf{Pesan utama:} Nilai utama model adalah ketika hasilnya dapat digunakan untuk mendukung keputusan nyata.\end{quote}


\hrule


\subsection{13. Kode Program}

\subsubsection{13.1 Generate Data Sintetis}

\begin{lstlisting}[language=Python]
import numpy as np
import pandas as pd

np.random.seed(42)
n = 180

rainfall = np.random.normal(180, 45, n)
temperature = np.random.normal(27, 2.2, n)
humidity = np.random.normal(75, 8, n)
fertilizer = np.random.normal(110, 25, n)
land_area = np.random.uniform(0.5, 3.0, n)

# Membentuk target dengan efek non-linear
rain_effect = -0.0009 * (rainfall - 180)**2 + 5.5
fert_effect = 0.025 * fertilizer
temp_effect = -0.08 * (temperature - 27)**2 + 1.2

yield_ton_ha = rain_effect + fert_effect + temp_effect
\end{lstlisting}

\subsubsection{13.2 Feature Engineering}

\begin{lstlisting}[language=Python]
df["rainfall_squared"] = df["rainfall"] ** 2

df["fertilizer_per_area"] = df["fertilizer"] / df["land_area"]

df["rainfall_category"] = pd.cut(
    df["rainfall"],
    bins=[-np.inf, 140, 220, np.inf],
    labels=["rendah", "sedang", "tinggi"]
)
\end{lstlisting}

\subsubsection{13.3 Supervised Learning}

\begin{lstlisting}[language=Python]
from sklearn.linear_model import LinearRegression
from sklearn.metrics import mean_absolute_error, mean_squared_error

features = [
    "rainfall",
    "rainfall_squared",
    "temperature",
    "humidity",
    "fertilizer",
    "land_area"
]

X = df[features]
y = df["yield_ton_ha"]

model = LinearRegression()
model.fit(X, y)

y_pred = model.predict(X)
\end{lstlisting}

\subsubsection{13.4 Unsupervised Learning}

\begin{lstlisting}[language=Python]
from sklearn.cluster import KMeans
from sklearn.preprocessing import StandardScaler
from sklearn.metrics import silhouette_score

X_cluster = df[["rainfall", "fertilizer", "temperature", "yield_ton_ha"]]

scaler = StandardScaler()
X_scaled = scaler.fit_transform(X_cluster)

kmeans = KMeans(n_clusters=3, random_state=42)
labels = kmeans.fit_predict(X_scaled)

score = silhouette_score(X_scaled, labels)
\end{lstlisting}

\hrule


\subsection{14. Ringkasan: Satu Use Case, Banyak Pendekatan}

\begin{tabular}{|l|l|l|}
\hline
Tujuan & Pendekatan & Output \\
\hline
Prediksi hasil panen & Supervised Learning & Nilai hasil panen \\
\hline
Segmentasi lahan & Unsupervised Learning & Kelompok lahan \\
\hline
Optimasi irigasi & Reinforcement Learning & Kebijakan aksi \\
\hline
\end{tabular}


\hrule


\subsection{15. Kesimpulan Utama}

\begin{enumerate}
\item EDA adalah fondasi penting sebelum modeling.
\end{enumerate}

\begin{enumerate}
\item Feature engineering harus didasarkan pada insight dari data.
\end{enumerate}

\begin{enumerate}
\item Supervised learning digunakan untuk prediksi dengan label.
\end{enumerate}

\begin{enumerate}
\item Unsupervised learning digunakan untuk menemukan struktur tersembunyi.
\end{enumerate}

\begin{enumerate}
\item Evaluasi model harus disesuaikan dengan jenis masalah.
\end{enumerate}


\begin{quote}
\textbf{Pesan akhir:} Model yang baik bukan hanya akurat, tetapi juga dapat dijelaskan dan berguna untuk pengambilan keputusan.\end{quote}



\end{document}
